\documentclass[12pt]{article}
\usepackage[utf8]{inputenc}
\usepackage{amsmath}
\usepackage{amsfonts}
\usepackage{amssymb}
\usepackage{geometry}
\usepackage{booktabs}
\usepackage{array}
\usepackage{float}

\geometry{margin=1in}

\title{Environmental Impact Mitigation Costs Methodology\\Comprehensive Transmission Cost Calculator}
\author{Andrew Igdal}
\date{\today}

\begin{document}

\maketitle

\section{Overview}
This document describes the methodology for calculating environmental impact mitigation costs for transmission line projects. The calculation considers different terrain types and their associated environmental mitigation requirements, providing a comprehensive assessment of environmental costs.

\section{Methodology}

\subsection{Terrain Types}
The environmental impact mitigation cost calculation considers nine distinct terrain types, each with different environmental sensitivities and mitigation requirements:

\begin{enumerate}
    \item \textbf{Forested}: Areas with significant tree cover requiring forest management and wildlife protection
    \item \textbf{Scrubbed/Flat}: Open areas with minimal vegetation requiring basic environmental monitoring
    \item \textbf{Wetland}: Wetland areas requiring extensive environmental protection and restoration
    \item \textbf{Farmland}: Agricultural areas requiring soil protection and agricultural impact mitigation
    \item \textbf{Desert/Barren Land}: Arid areas requiring minimal environmental mitigation
    \item \textbf{Urban}: Urban areas requiring community impact mitigation and aesthetic considerations
    \item \textbf{Rolling Hills (2-8\% Slope)}: Moderately sloped areas requiring erosion control
    \item \textbf{Mountain (>8\% Slope)}: Steep terrain requiring extensive slope stabilization and erosion control
    \item \textbf{Subsea}: Underwater transmission requiring marine environmental protection
\end{enumerate}

\subsection{Mathematical Formulation}

The total environmental mitigation cost for each terrain type is calculated as:
\begin{equation}
C_{terrain,i} = L_i \times c_{terrain,i}
\end{equation}

Where:
\begin{itemize}
    \item $L_i$ = Length of transmission line in terrain type $i$ (miles)
    \item $c_{terrain,i}$ = Environmental mitigation cost per mile for terrain type $i$ (USD/mile)
\end{itemize}

The total environmental mitigation cost is:
\begin{equation}
C_{total} = \sum_{i=1}^{9} C_{terrain,i}
\end{equation}

The total project length is:
\begin{equation}
L_{total} = \sum_{i=1}^{9} L_i
\end{equation}

The present value of total environmental mitigation costs over the project lifetime is:
\begin{equation}
PV_{environmental} = \sum_{t=1}^{T} \frac{C_{total}}{(1 + r)^t}
\end{equation}

Where:
\begin{itemize}
    \item $T$ = Project lifetime in years
    \item $r$ = Firm interest rate (decimal)
    \item $t$ = Time period (year)
\end{itemize}

\section{Input Parameters}

\begin{table}[H]
\centering
\caption{Required Input Parameters}
\begin{tabular}{@{}ll@{}}
\toprule
\textbf{Parameter} & \textbf{Description} \\
\midrule
Firm Interest Rate & Discount rate for present value calculation (decimal) \\
Project Lifetime & Expected project lifetime in years \\
Forested Miles & Length of transmission line in forested terrain (miles) \\
Scrubbed/Flat Miles & Length of transmission line in scrubbed/flat terrain (miles) \\
Wetland Miles & Length of transmission line in wetland terrain (miles) \\
Farmland Miles & Length of transmission line in farmland terrain (miles) \\
Desert/Barren Miles & Length of transmission line in desert/barren terrain (miles) \\
Urban Miles & Length of transmission line in urban terrain (miles) \\
Rolling Hills Miles & Length of transmission line in rolling hills terrain (miles) \\
Mountain Miles & Length of transmission line in mountain terrain (miles) \\
Subsea Miles & Length of transmission line in subsea terrain (miles) \\
\bottomrule
\end{tabular}
\end{table}

\begin{table}[H]
\centering
\caption{Environmental Mitigation Cost Parameters}
\begin{tabular}{@{}ll@{}}
\toprule
\textbf{Parameter} & \textbf{Description} \\
\midrule
Forested Cost/Mile & Environmental mitigation cost per mile for forested terrain (USD) \\
Scrubbed/Flat Cost/Mile & Environmental mitigation cost per mile for scrubbed/flat terrain (USD) \\
Wetland Cost/Mile & Environmental mitigation cost per mile for wetland terrain (USD) \\
Farmland Cost/Mile & Environmental mitigation cost per mile for farmland terrain (USD) \\
Desert/Barren Cost/Mile & Environmental mitigation cost per mile for desert/barren terrain (USD) \\
Urban Cost/Mile & Environmental mitigation cost per mile for urban terrain (USD) \\
Rolling Hills Cost/Mile & Environmental mitigation cost per mile for rolling hills terrain (USD) \\
Mountain Cost/Mile & Environmental mitigation cost per mile for mountain terrain (USD) \\
Subsea Cost/Mile & Environmental mitigation cost per mile for subsea terrain (USD) \\
\bottomrule
\end{tabular}
\end{table}

\section{Outputs}

The script provides the following outputs:
\begin{itemize}
    \item Detailed breakdown of costs by terrain type
    \item Total project miles
    \item Total environmental mitigation cost
    \item Present value of total environmental mitigation cost
\end{itemize}

\section{Environmental Considerations by Terrain Type}

\begin{itemize}
    \item \textbf{Forested}: Wildlife habitat protection, tree clearing permits, reforestation requirements
    \item \textbf{Scrubbed/Flat}: Minimal environmental impact, basic monitoring requirements
    \item \textbf{Wetland}: Wetland delineation, mitigation banking, water quality protection
    \item \textbf{Farmland}: Agricultural impact assessment, soil protection, crop compensation
    \item \textbf{Desert/Barren}: Minimal vegetation impact, desert tortoise protection (if applicable)
    \item \textbf{Urban}: Community impact mitigation, aesthetic considerations, noise abatement
    \item \textbf{Rolling Hills}: Erosion control, slope stabilization, runoff management
    \item \textbf{Mountain}: Extensive slope stabilization, rockfall protection, access road environmental impact
    \item \textbf{Subsea}: Marine environmental protection, benthic habitat impact, fishing impact mitigation
\end{itemize}

\section{Assumptions}

\begin{itemize}
    \item Environmental mitigation costs are incurred annually throughout the project lifetime
    \item The firm interest rate remains constant over the project lifetime
    \item All costs are expressed in nominal USD
    \item Environmental mitigation costs are additive across all terrain types
    \item Cost per mile remains constant for each terrain type
\end{itemize}

\section{Usage Notes}

\begin{itemize}
    \item Input all cost values in USD per mile
    \item Interest rate should be entered as a decimal (e.g., 0.05 for 5\%)
    \item Project lifetime should be entered in years
    \item All mile values should be entered in miles
    \item Ensure that the sum of all terrain miles equals the total project length
\end{itemize}

\section{Example Calculation}

For a 200-mile transmission project with the following terrain distribution:
\begin{itemize}
    \item Forested: 50 miles at \$5,000/mile
    \item Urban: 30 miles at \$8,000/mile
    \item Farmland: 120 miles at \$3,000/mile
\end{itemize}

The environmental mitigation costs would be:
\begin{align}
C_{forested} &= 50 \times 5,000 = \$250,000 \\
C_{urban} &= 30 \times 8,000 = \$240,000 \\
C_{farmland} &= 120 \times 3,000 = \$360,000 \\
C_{total} &= 250,000 + 240,000 + 360,000 = \$850,000
\end{align}

\end{document}
